%-------------------------------------------------------------------------------
%	SECTION TITLE
%-------------------------------------------------------------------------------
\cvsection{Projeler}

%-------------------------------------------------------------------------------
%	CONTENT
%-------------------------------------------------------------------------------
\begin{cventries}

%---------------------------------------------------------
\cventry
  {İGA İstanbul Havalimanı}
  {F1RST – Gerçek Zamanlı Operasyonel Karar Destek Platformu}
  {İstanbul}
  {May 2024 - Halen}
  {
    \begin{cvitems}
      \item {MERN Stack kullanarak, gerçek zamanlı uçak durumu ve planlamayı görselleştiren özel UI bileşenleri ile kritik bir platformun geliştirilmesi.}
      \item {Canlı AMS akışlarını işleyen, saniyenin altında gecikme sağlayan Apache Kafka ve Spark tabanlı yüksek verimli bir backend mimarisinin tasarlanması.}
      \item {Eski gecikmeli raporlama sistemlerinin anlık veri görselleştirme ile değiştirilerek operasyonel ekiplere veri erişiminin kolaylaştırılması.}
    \end{cvitems}
  }
%---------------------------------------------------------
    \cventry
    {İGA İstanbul Havalimanı}
    {Birleşik Gerçek Zamanlı \& Toplu İş Entegrasyon Çerçevesi}
    {İstanbul}
    {Haz 2023 – Halen}
    {
      \begin{cvitems}
        \item {ASMGCS, AMS, SAP ve MSSQL gibi farklı veri kaynaklarını merkezi bir ortamda senkronize eden sağlam bir çerçevenin tasarlanması.}
        \item {Kurumsal raporlama için "tek doğru kaynak" sağlayan gerçek zamanlı ve toplu iş yükleri için veri alım katmanlarının geliştirilmesi.}
        \item {Redis ve Spark aracılığıyla veri hareketinin optimize edilerek sağlam ve yüksek performanslı analitik çözümler sunulması.}
        \item {Eski veri hatlarının gerçek zamanlı AMS verilerini işleyecek şekilde modernize edilerek, veri erişim gecikmesinin 15 dakikadan saniyenin altına düşürülmesi.}
      \end{cvitems}
    }

%---------------------------------------------------------
  \cventry
    {İGA İstanbul Havalimanı}
    {Kurumsal Veri Ambarı \& ETL Mimarisi}
    {İstanbul}
    {Haz 2023 – Halen}
    {
      \begin{cvitems}
        \item {Staging (STG), Operasyonel Veri Deposu (ODS) ve Veri Ambarı (DWH) katmanlarından oluşan çok katmanlı bir veri mimarisinin mühendisliği.}
        \item {Üst düzey yönetim analitiği ve operasyonel raporlamayı mümkün kılan "tek doğru kaynak" oluşturmak için karmaşık dönüşüm mantığının uygulanması.}
        \item {PostgreSQL ve MongoDB depoları genelinde \%100 veri bütünlüğü sağlamak için her katmanda veri kalitesi protokollerinin oluşturulması.}
      \end{cvitems}
    }

\end{cventries}
